\documentclass[11pt,oneside]{book}

% ============================================================
% Packages
% ============================================================
\usepackage[a4paper,margin=1in]{geometry}
\usepackage[T1]{fontenc}
\usepackage[utf8]{inputenc}
\usepackage{lmodern}
\usepackage{microtype}
\usepackage{amsmath,amssymb,amsthm,mathtools}
\usepackage{bm}
\usepackage{enumitem}
\usepackage{booktabs,tabularx,array,multicol}
\usepackage{xcolor}
\usepackage{tikz}
\usetikzlibrary{arrows.meta,calc,patterns,positioning}
\usepackage{pgfplots}
\pgfplotsset{compat=1.18}
\usepackage{titlesec}
\usepackage{fancyhdr}
\usepackage{hyperref}
\usepackage{setspace}
\usepackage{needspace}

% ============================================================
% Colours and typography
% ============================================================
\definecolor{MainBlue}{HTML}{163A5F}
\definecolor{DeepBlue}{HTML}{0D2B45}
\definecolor{MainGreen}{HTML}{2F6B4F}
\definecolor{MainGold}{HTML}{9A6A16}
\definecolor{MainRed}{HTML}{8A3030}
\definecolor{Ink}{HTML}{222222}

\color{Ink}
\setstretch{1.08}
\setlength{\parindent}{0pt}
\setlength{\parskip}{4pt}
\setlist[itemize]{leftmargin=1.6em,itemsep=2pt,topsep=3pt}
\setlist[enumerate]{leftmargin=2em,itemsep=3pt,topsep=3pt}

% ============================================================
% Page style
% ============================================================
\pagestyle{fancy}
\fancyhf{}
\fancyhead[L]{\small\textsc{Undergraduate Calculus Notes}}
\fancyhead[R]{\small\thepage}
\fancyfoot{}
\renewcommand{\headrulewidth}{0.4pt}
\setlength{\headheight}{14pt}

\titleformat{\chapter}[display]
  {\normalfont\bfseries\color{DeepBlue}}
  {\filleft\Large\chaptername\ \thechapter}
  {1ex}
  {\titlerule\vspace{1.2ex}\Huge}
  [\vspace{1ex}\titlerule]

\titleformat{\section}
  {\Large\bfseries\color{MainBlue}}
  {\thesection}{0.7em}{}

\titleformat{\subsection}
  {\large\bfseries\color{MainGreen}}
  {\thesubsection}{0.7em}{}

% ============================================================
% Theorem environments
% ============================================================
\newtheoremstyle{studentthm}
{6pt}{6pt}{\itshape}{}
{\bfseries\color{DeepBlue}}{.}{0.5em}{}
\theoremstyle{studentthm}
\newtheorem{theorem}{Theorem}[chapter]
\newtheorem{proposition}[theorem]{Proposition}
\newtheorem{lemma}[theorem]{Lemma}
\newtheorem{corollary}[theorem]{Corollary}

\theoremstyle{definition}
\newtheorem{definition}[theorem]{Definition}
\newtheorem{example}[theorem]{Example}
\newtheorem{remark}[theorem]{Remark}

% ============================================================
% Commands
% ============================================================
\newcommand{\R}{\mathbb{R}}
\newcommand{\N}{\mathbb{N}}
\newcommand{\eps}{\varepsilon}
\newcommand{\abs}[1]{\left|#1\right|}
\newcommand{\norm}[1]{\left\lVert#1\right\rVert}
\newcommand{\pd}[2]{\frac{\partial #1}{\partial #2}}
\newcommand{\pdd}[3]{\frac{\partial^2 #1}{\partial #2\,\partial #3}}
\newcommand{\dd}{\,\mathrm d}

\hypersetup{
  colorlinks=true,
  linkcolor=MainBlue,
  urlcolor=MainGreen,
  pdfauthor={C. Zosangzuala},
  pdftitle={Functions of Two Variables: Limits, Partial Derivatives, Differentiability and Homogeneous Functions}
}

% ============================================================
\begin{document}

% ============================================================
% Cover
% ============================================================
\begin{titlepage}
\centering
\vspace*{2.0cm}

{\Huge\bfseries\color{DeepBlue}
Functions of Two Variables\par}

\vspace{0.35cm}

{\LARGE
Limits, Continuity, Partial Derivatives,\\
Differentiability, Composite Functions,\\
Change of Variables and Homogeneous Functions\par}

\vspace{1.25cm}

{\large
Extensive Lecture Notes for Undergraduate Calculus\par}

\vspace{0.9cm}

{\large
Real-valued functions of two variables; limits and repeated limits;
continuity; first- and second-order partial derivatives;
differentiability; Young's and Schwarz's theorems;
derivatives of composite functions; change of variables;
Jacobians; and Euler's theorem on homogeneous functions.\par}

\vfill

{\Large\bfseries C. Zosangzuala\par}
\vspace{0.25cm}
{\large Lecture Notes and Problem-Solving Companion\par}

\vspace{1cm}
\end{titlepage}

\frontmatter
\tableofcontents

% ============================================================
\chapter*{How to Use These Notes}
\addcontentsline{toc}{chapter}{How to Use These Notes}

Calculus of one variable studies a function such as
\[
y=f(x),
\]
where one real number determines the output. In this unit the input
consists of two real numbers:
\[
z=f(x,y).
\]
This apparently small change introduces several new ideas.

For a function of one variable, \(x\) approaches \(a\) essentially
from the left or from the right. In two variables, however,
\((x,y)\) can approach \((a,b)\) along infinitely many lines,
parabolas and other curves. Consequently, two-variable limits require
greater care.

Similarly, the derivative of a one-variable function gives one local
slope, whereas a surface \(z=f(x,y)\) has different rates of change
depending on how the point \((x,y)\) moves.

The conceptual order of the unit is
\[
\text{functions}
\longrightarrow
\text{limits and continuity}
\longrightarrow
\text{partial derivatives}
\longrightarrow
\text{differentiability}
\]
\[
\longrightarrow
\text{composite functions}
\longrightarrow
\text{change of variables}
\longrightarrow
\text{homogeneous functions}.
\]

The most important distinction to remember is
\[
\text{existence of partial derivatives}
\quad\not\Rightarrow\quad
\text{differentiability}.
\]

\mainmatter

% ============================================================
\chapter{Real-Valued Functions of Two Variables}
% ============================================================

\section*{Learning Objectives}

After this chapter a student should be able to:
\begin{itemize}
\item define a real-valued function of two variables;
\item find and sketch common domains in the \(xy\)-plane;
\item interpret the graph \(z=f(x,y)\) as a surface;
\item use level curves to understand a surface;
\item work with distance, neighbourhoods and deleted neighbourhoods
in \(\R^2\).
\end{itemize}

\section*{Historical Motivation}

The extension of calculus to several variables grew out of problems in
mechanics, astronomy, geometry and mathematical physics. Temperature,
pressure, potential energy and density naturally depend on more than
one independent variable. Euler, Lagrange, Clairaut, Jacobi and
Schwarz helped develop the language now used in multivariable
calculus.

\section{Definition}

\begin{definition}
Let \(D\subseteq\R^2\). A \emph{real-valued function of two real
variables} is a rule
\[
f:D\longrightarrow\R
\]
which assigns exactly one real number \(f(x,y)\) to each
\((x,y)\in D\).
\end{definition}

We usually write
\[
z=f(x,y).
\]

Here \(x,y\) are independent variables and \(z\) is the dependent
variable.

\paragraph{Everyday comparison.}
A quantity may depend on two independent inputs. For example, a taxi
fare may depend on both distance travelled and waiting time:
\[
\text{fare}=f(\text{distance},\text{waiting time}).
\]
Two inputs determine one output.

\begin{example}
For
\[
f(x,y)=x^2+y^2,
\]
every ordered pair is allowed, so
\[
\operatorname{Dom}(f)=\R^2.
\]
Since \(x^2+y^2\ge0\),
\[
\operatorname{Range}(f)=[0,\infty).
\]
\end{example}

\section{Finding Domains}

The familiar restrictions from one-variable calculus still apply:
\begin{itemize}
\item a denominator cannot be zero;
\item an even root requires a non-negative radicand;
\item a logarithm requires a positive argument.
\end{itemize}

\begin{example}
Find the domain of
\[
f(x,y)=\sqrt{9-x^2-y^2}.
\]
We require
\[
9-x^2-y^2\ge0,
\]
hence
\[
x^2+y^2\le9.
\]
The domain is the closed disc of radius \(3\) centred at the origin.
\end{example}

\begin{example}
Find the domain of
\[
f(x,y)=\log(4-x^2-y^2).
\]
We require
\[
4-x^2-y^2>0,
\]
hence
\[
x^2+y^2<4.
\]
The boundary circle is excluded.
\end{example}

\begin{example}
Find the domain of
\[
f(x,y)=\frac{x+y}{x-y}.
\]
The denominator must be nonzero:
\[
x-y\ne0.
\]
Thus the line \(y=x\) is removed from the plane.
\end{example}

\section{Graphs and Surfaces}

The graph of \(z=f(x,y)\) is
\[
\{(x,y,z)\in\R^3:z=f(x,y)\}.
\]
It is generally a surface.

For example,
\[
z=x^2+y^2
\]
is an upward-opening paraboloid.

\begin{center}
\begin{tikzpicture}
\begin{axis}[
width=10cm,height=7cm,
view={45}{27},
xlabel={$x$},ylabel={$y$},zlabel={$z$},
domain=-2:2,y domain=-2:2,
samples=23,samples y=23
]
\addplot3[surf,shader=interp,opacity=0.85]{x^2+y^2};
\end{axis}
\end{tikzpicture}
\end{center}

\section{Level Curves}

\begin{definition}
For a function \(z=f(x,y)\), a \emph{level curve} corresponding to
\(c\) is defined by
\[
f(x,y)=c.
\]
\end{definition}

For
\[
f(x,y)=x^2+y^2,
\]
the level curves are
\[
x^2+y^2=c.
\]
For \(c>0\), these are concentric circles.

\begin{center}
\begin{tikzpicture}[scale=0.9]
\draw[->] (-3.4,0)--(3.4,0) node[right] {$x$};
\draw[->] (0,-3.4)--(0,3.4) node[above] {$y$};
\draw[thick] (0,0) circle (1);
\draw[thick] (0,0) circle (1.8);
\draw[thick] (0,0) circle (2.6);
\node at (0.7,0.7) {$c=1$};
\node at (1.4,1.35) {$c=3.24$};
\node at (2.15,2.05) {$c=6.76$};
\end{tikzpicture}
\end{center}

Level curves are analogous to contour lines on a geographical map.

\section{Distance and Neighbourhoods}

The distance from \(P=(x,y)\) to \(P_0=(a,b)\) is
\[
d(P,P_0)
=
\sqrt{(x-a)^2+(y-b)^2}.
\]

\begin{definition}
The \(\delta\)-neighbourhood of \((a,b)\) is
\[
N_\delta(a,b)
=
\left\{
(x,y):
\sqrt{(x-a)^2+(y-b)^2}<\delta
\right\}.
\]
\end{definition}

A deleted neighbourhood excludes the centre:
\[
0<
\sqrt{(x-a)^2+(y-b)^2}
<\delta.
\]

\paragraph{Neighbourhood analogy.}
Think of \((a,b)\) as a particular house. A neighbourhood contains
all points sufficiently close to the house. In the plane there are
infinitely many possible directions from which we may approach it.

\section{Detailed Solved Problems}

\subsection*{Problem 1. Domain of a square-root function}
Find the domain of
\[
f(x,y)=\sqrt{16-x^2-4y^2}.
\]

We require
\[
16-x^2-4y^2\ge0.
\]
Thus
\[
x^2+4y^2\le16.
\]
Dividing by \(16\),
\[
\frac{x^2}{16}+\frac{y^2}{4}\le1.
\]
The domain is the closed elliptical region inside and on the ellipse.

\subsection*{Problem 2. Domain of a logarithmic function}
Find the domain of
\[
f(x,y)=\log(x-y^2).
\]
We require
\[
x-y^2>0,
\]
so
\[
x>y^2.
\]
The domain is the region to the right of the parabola \(x=y^2\);
the parabola itself is excluded.

\subsection*{Problem 3. Level curves}
Describe the level curves of
\[
f(x,y)=x^2+4y^2.
\]
Setting \(f(x,y)=c\),
\[
x^2+4y^2=c.
\]
For \(c>0\), these are ellipses centred at the origin.

\section{Graded Practice Problems}

\subsection*{Level I: Foundations}
\begin{enumerate}
\item Find the domain of \(x^2-3xy+y^2\).
\item Find the domain of \(\sqrt{25-x^2-y^2}\).
\item Find the domain of \(1/(x+y)\).
\item Describe the level curves of \(f(x,y)=x+y\).
\end{enumerate}

\subsection*{Level II: Standard Problems}
\begin{enumerate}
\item Find and sketch the domain of \(\sqrt{y-x^2}\).
\item Find the domain of \(\log(9-x^2-y^2)\).
\item Describe the level curves of \(4x^2+y^2\).
\item Explain geometrically the surface \(z=4-x^2-y^2\).
\end{enumerate}

\subsection*{Level III: Challenging}
\begin{enumerate}
\item Determine the domain of
\[
\sqrt{\frac{x-y}{x+y}}.
\]
\item Describe the level curves of \(x/y\).
\item Compare the surfaces \(z=x^2+y^2\) and \(z=x^2-y^2\).
\end{enumerate}

% ============================================================
\chapter{Limits, Repeated Limits and Continuity}
% ============================================================

\section*{Learning Objectives}
Students should be able to:
\begin{itemize}
\item state and use the \(\eps\)-\(\delta\) definition of a
two-variable limit;
\item test nonexistence by different paths;
\item understand why a few agreeing paths do not prove existence;
\item use inequalities and polar coordinates to prove existence;
\item compute repeated limits;
\item distinguish repeated limits from a joint limit;
\item test continuity.
\end{itemize}

\section*{Historical Motivation}

The rigorous ideas of Cauchy and Weierstrass extend naturally to the
plane, but the geometry changes. A point in the plane can be
approached in infinitely many ways. A correct definition must
therefore control all approaches simultaneously.

\section{Definition of a Joint Limit}

\begin{definition}
We say
\[
\lim_{(x,y)\to(a,b)}f(x,y)=L
\]
if, for every \(\eps>0\), there exists \(\delta>0\) such that
\[
0<
\sqrt{(x-a)^2+(y-b)^2}
<\delta
\]
implies
\[
\abs{f(x,y)-L}<\eps.
\]
\end{definition}

The distance condition does not prescribe a path, so the definition
controls every possible approach to \((a,b)\).

\section{Visual Meaning of \(\eps\) and \(\delta\)}

\begin{center}
\begin{tikzpicture}[scale=1.0]
\draw[->] (-0.5,0)--(6,0) node[right] {$x$};
\draw[->] (0,-0.5)--(0,5) node[above] {$y$};
\coordinate (A) at (3,2.2);
\draw[thick] (A) circle (1.35);
\fill (A) circle (2pt);
\node[below right] at (A) {$(a,b)$};
\draw[<->] (3,2.2)--(4.35,2.2);
\node[above] at (3.7,2.2) {$\delta$};
\node at (3,4.1) {$N_\delta(a,b)$};
\end{tikzpicture}
\end{center}

The \(\delta\)-disc controls the inputs. The \(\eps\)-condition controls
the outputs.

\section{A Direct \(\eps\)-\(\delta\) Proof}

\begin{example}
Prove
\[
\lim_{(x,y)\to(0,0)}(3x-2y)=0.
\]
\end{example}

\textbf{Proof.}
We want to make
\[
\abs{3x-2y}
\]
smaller than a prescribed \(\eps>0\).

By the triangle inequality,
\[
\abs{3x-2y}
\le3\abs{x}+2\abs{y}.
\]

Also,
\[
\abs{x}\le\sqrt{x^2+y^2},
\qquad
\abs{y}\le\sqrt{x^2+y^2}.
\]
Hence
\[
\abs{3x-2y}
\le5\sqrt{x^2+y^2}.
\]

Choose
\[
\delta=\frac{\eps}{5}.
\]
Then
\[
0<\sqrt{x^2+y^2}<\delta
\]
implies
\[
\abs{3x-2y}
<
5\delta
=
\eps.
\]
Therefore
\[
\lim_{(x,y)\to(0,0)}(3x-2y)=0.
\]

\section{Standard Limit Laws}

Suppose
\[
\lim_{(x,y)\to(a,b)}f(x,y)=L
\]
and
\[
\lim_{(x,y)\to(a,b)}g(x,y)=M.
\]
Then:
\[
\lim(f+g)=L+M,
\]
\[
\lim(f-g)=L-M,
\]
\[
\lim(fg)=LM,
\]
and, if \(M\ne0\),
\[
\lim\frac{f}{g}=\frac{L}{M}.
\]

The proofs follow from the same \(\eps\)-\(\delta\) ideas used in
one-variable calculus.

\section{Path Testing}

If the joint limit exists, every path to the point must give the same
value. Thus two paths with different limiting values prove
nonexistence immediately.

\begin{example}
Examine
\[
\lim_{(x,y)\to(0,0)}
\frac{x^2-y^2}{x^2+y^2}.
\]

Along \(y=0\),
\[
\frac{x^2}{x^2}=1.
\]
Along \(x=0\),
\[
\frac{-y^2}{y^2}=-1.
\]
The two limits are different. Therefore the joint limit does not
exist.
\end{example}

\section{Why Agreement Along Straight Lines Is Not Enough}

Consider
\[
f(x,y)
=
\frac{x^2y}{x^4+y^2}.
\]
Along the line \(y=mx\),
\[
f(x,mx)
=
\frac{mx^3}{x^4+m^2x^2}
=
\frac{mx}{x^2+m^2}
\longrightarrow0.
\]
Thus every fixed straight line gives \(0\).

But along the parabola \(y=x^2\),
\[
f(x,x^2)
=
\frac{x^4}{x^4+x^4}
=
\frac12.
\]
Hence the joint limit does not exist.

\paragraph{Key principle.}
Different paths are excellent for proving nonexistence. Agreement
along selected paths is not, by itself, a proof of existence.

\section{Proving Existence by an Inequality}

Consider
\[
f(x,y)=\frac{x^2y}{x^2+y^2}.
\]
Then
\[
\abs{f(x,y)}
=
\abs{y}\frac{x^2}{x^2+y^2}
\le\abs{y}.
\]
Since
\[
\abs{y}\to0
\quad\text{as}\quad
(x,y)\to(0,0),
\]
the Squeeze (Sandwich) Theorem gives
\[
\lim_{(x,y)\to(0,0)}
\frac{x^2y}{x^2+y^2}
=0.
\]

\section{Polar Coordinates in Limit Problems}

At the origin, set
\[
x=r\cos\theta,\qquad y=r\sin\theta.
\]
Then
\[
x^2+y^2=r^2
\]
and
\[
(x,y)\to(0,0)
\quad\Longleftrightarrow\quad
r\to0.
\]

For example,
\[
\frac{x^2y}{x^2+y^2}
=
r\cos^2\theta\sin\theta.
\]
Hence
\[
\abs{r\cos^2\theta\sin\theta}\le r,
\]
so the limit is \(0\).

\section{Repeated Limits}

\begin{definition}
The repeated limit
\[
\lim_{x\to a}
\left[
\lim_{y\to b}f(x,y)
\right]
\]
is obtained by first holding \(x\) fixed and taking \(y\to b\), then
taking \(x\to a\).
\end{definition}

The reverse order is
\[
\lim_{y\to b}
\left[
\lim_{x\to a}f(x,y)
\right].
\]

\begin{example}
For
\[
f(x,y)=\frac{x^2-y^2}{x^2+y^2},
\]
fix \(x\ne0\). Then
\[
\lim_{y\to0}f(x,y)=1.
\]
Thus
\[
\lim_{x\to0}\lim_{y\to0}f(x,y)=1.
\]

On the other hand, fixing \(y\ne0\),
\[
\lim_{x\to0}f(x,y)=-1,
\]
so
\[
\lim_{y\to0}\lim_{x\to0}f(x,y)=-1.
\]
The repeated limits are different.
\end{example}

\section{Relation Between Joint and Repeated Limits}

\begin{theorem}
If the joint limit
\[
\lim_{(x,y)\to(a,b)}f(x,y)=L
\]
exists and the relevant inner limits exist in a neighbourhood, then
both repeated limits, when defined, must equal \(L\).
\end{theorem}

However, the converse is false.

\begin{example}
Let
\[
f(x,y)=
\begin{cases}
\dfrac{xy}{x^2+y^2},&(x,y)\ne(0,0),\\[1ex]
0,&(x,y)=(0,0).
\end{cases}
\]
For fixed \(x\),
\[
\lim_{y\to0}f(x,y)=0,
\]
so
\[
\lim_{x\to0}\lim_{y\to0}f(x,y)=0.
\]
Similarly,
\[
\lim_{y\to0}\lim_{x\to0}f(x,y)=0.
\]

But along \(y=x\),
\[
f(x,x)=\frac12.
\]
Therefore the joint limit does not exist.

Thus
\[
\text{equal repeated limits}
\not\Rightarrow
\text{existence of the joint limit}.
\]
\end{example}

\section{Continuity}

\begin{definition}
A function \(f(x,y)\) is continuous at \((a,b)\) if:
\begin{enumerate}
\item \(f(a,b)\) is defined;
\item \(\lim_{(x,y)\to(a,b)}f(x,y)\) exists;
\item
\[
\lim_{(x,y)\to(a,b)}f(x,y)=f(a,b).
\]
\end{enumerate}
\end{definition}

Polynomials in \(x,y\) are continuous everywhere. Rational functions
are continuous wherever their denominators are nonzero.

\begin{example}
Define
\[
f(x,y)=
\begin{cases}
\dfrac{x^2y}{x^2+y^2},&(x,y)\ne(0,0),\\[1ex]
0,&(x,y)=(0,0).
\end{cases}
\]
Since
\[
\abs{\frac{x^2y}{x^2+y^2}}
\le\abs{y}\to0,
\]
we have
\[
\lim_{(x,y)\to(0,0)}f(x,y)=0=f(0,0).
\]
Hence \(f\) is continuous at the origin.
\end{example}

\section{Detailed Solved Problems}

\subsection*{Problem 1. A direct limit}
Evaluate
\[
\lim_{(x,y)\to(1,2)}(x^2+xy+y^2).
\]
Because a polynomial is continuous,
\[
1^2+(1)(2)+2^2=7.
\]

\subsection*{Problem 2. Limit by comparison}
Evaluate
\[
\lim_{(x,y)\to(0,0)}
\frac{x^3}{x^2+y^2}.
\]
Since
\[
\abs{\frac{x^3}{x^2+y^2}}
=
\abs{x}\frac{x^2}{x^2+y^2}
\le\abs{x},
\]
and \(\abs{x}\to0\), the limit is \(0\).

\subsection*{Problem 3 (2024). Nonexistence of a double limit}
Examine
\[
\lim_{(x,y)\to(0,0)}
\frac{2x}{x^2+y^2}.
\]
Along \(x=0\), the expression equals \(0\). Along \(y=0\),
\[
\frac{2x}{x^2}=\frac2x,
\]
which has no finite limit as \(x\to0\). Therefore the joint limit
does not exist.

\subsection*{Problem 4. Repeated limits}
Let
\[
f(x,y)=\frac{x-y}{x+y}.
\]
For fixed \(x\ne0\),
\[
\lim_{y\to0}f(x,y)=1,
\]
hence
\[
\lim_{x\to0}\lim_{y\to0}f(x,y)=1.
\]
For fixed \(y\ne0\),
\[
\lim_{x\to0}f(x,y)=-1,
\]
hence
\[
\lim_{y\to0}\lim_{x\to0}f(x,y)=-1.
\]

\section{Graded Practice Problems}

\subsection*{Level I}
\begin{enumerate}
\item Evaluate
\[
\lim_{(x,y)\to(1,1)}(x^2+3xy-y^2).
\]
\item Evaluate
\[
\lim_{(x,y)\to(0,0)}
\frac{x^2y}{x^2+y^2}.
\]
\item Test
\[
\lim_{(x,y)\to(0,0)}
\frac{x^2-y^2}{x^2+y^2}.
\]
\item State the definition of continuity at \((a,b)\).
\end{enumerate}

\subsection*{Level II}
\begin{enumerate}
\item Examine
\[
\lim_{(x,y)\to(0,0)}
\frac{xy}{x^2+y^2}.
\]
\item Prove by the \(\eps\)-\(\delta\) definition that
\[
\lim_{(x,y)\to(0,0)}(x+y)=0.
\]
\item Evaluate both repeated limits of
\[
\frac{x^2-y^2}{x^2+y^2}.
\]
\end{enumerate}

\subsection*{Level III}
\begin{enumerate}
\item Show that straight-line checking is insufficient for
\[
\frac{x^2y}{x^4+y^2}.
\]
\item Find an example where both repeated limits are equal but the
joint limit fails to exist.
\item Use polar coordinates to prove
\[
\lim_{(x,y)\to(0,0)}
\frac{x^4y^2}{x^4+y^2}=0.
\]
\end{enumerate}

% ============================================================
\chapter{Partial Derivatives of First and Second Orders}
% ============================================================

\section*{Learning Objectives}
Students should be able to:
\begin{itemize}
\item define partial derivatives from first principles;
\item interpret them geometrically as slopes of traces;
\item calculate first and second partial derivatives;
\item distinguish pure and mixed derivatives;
\item understand what the existence of partial derivatives does and
does not imply.
\end{itemize}

\section*{Historical Motivation}

A function \(z=f(x,y)\) may respond differently to changes in \(x\)
and changes in \(y\). Partial differentiation isolates one direction
at a time by temporarily fixing the other variable. This became
fundamental in mechanics, thermodynamics and differential equations.

\section{First-Order Partial Derivatives}

\begin{definition}
The partial derivative of \(f\) with respect to \(x\) at \((a,b)\) is
\[
f_x(a,b)
=
\lim_{h\to0}
\frac{f(a+h,b)-f(a,b)}{h},
\]
provided the limit exists.
\end{definition}

Similarly,
\[
f_y(a,b)
=
\lim_{k\to0}
\frac{f(a,b+k)-f(a,b)}{k}.
\]

When differentiating with respect to \(x\), treat \(y\) as a constant.
When differentiating with respect to \(y\), treat \(x\) as a constant.

\paragraph{Real-life comparison.}
If \(T(x,y)\) is the temperature at position \((x,y)\), then \(T_x\)
measures how temperature changes when moving in the \(x\)-direction
while keeping \(y\) fixed; \(T_y\) does the analogous job in the
\(y\)-direction.

\section{Geometrical Interpretation}

For the surface
\[
z=f(x,y),
\]
fix \(y=b\). Then
\[
z=f(x,b)
\]
is a curve in the vertical plane \(y=b\), and
\[
f_x(a,b)
\]
is the tangent slope of this trace at \(x=a\).

Similarly \(f_y(a,b)\) is the tangent slope of the trace obtained by
fixing \(x=a\).

\section{Calculation from the Definition}

Let
\[
f(x,y)=x^2+y^2.
\]
Then
\[
\begin{aligned}
f_x(a,b)
&=
\lim_{h\to0}
\frac{(a+h)^2+b^2-(a^2+b^2)}{h}\\
&=
\lim_{h\to0}
\frac{2ah+h^2}{h}\\
&=
\lim_{h\to0}(2a+h)\\
&=2a.
\end{aligned}
\]
Similarly,
\[
f_y(a,b)=2b.
\]

\section{Second-Order Partial Derivatives}

The four second-order partial derivatives are
\[
f_{xx},\qquad f_{xy},\qquad f_{yx},\qquad f_{yy}.
\]

Here
\[
f_{xy}
=
\frac{\partial}{\partial y}(f_x),
\qquad
f_{yx}
=
\frac{\partial}{\partial x}(f_y).
\]

\begin{example}
Let
\[
f(x,y)=x^3y^2+e^{xy}.
\]
Then
\[
f_x=3x^2y^2+ye^{xy},
\]
\[
f_y=2x^3y+xe^{xy}.
\]
Differentiate again:
\[
f_{xx}=6xy^2+y^2e^{xy},
\]
\[
f_{xy}=6x^2y+e^{xy}+xye^{xy},
\]
\[
f_{yx}=6x^2y+e^{xy}+xye^{xy},
\]
and
\[
f_{yy}=2x^3+x^2e^{xy}.
\]
Thus in this example
\[
f_{xy}=f_{yx}.
\]
\end{example}

\section{Partial Derivatives Need Not Imply Continuity}

Define
\[
f(x,y)=
\begin{cases}
\dfrac{xy}{x^2+y^2},&(x,y)\ne(0,0),\\[1ex]
0,&(x,y)=(0,0).
\end{cases}
\]

By definition,
\[
f_x(0,0)
=
\lim_{h\to0}
\frac{f(h,0)-f(0,0)}{h}
=0.
\]
Likewise,
\[
f_y(0,0)=0.
\]

But along \(y=x\),
\[
f(x,x)=\frac12,
\]
so the function is not continuous at \((0,0)\).

Therefore
\[
f_x(0,0),f_y(0,0)\text{ exist}
\quad\not\Rightarrow\quad
f\text{ is continuous at }(0,0).
\]

\section{Detailed Solved Problems}

\subsection*{Problem 1}
For
\[
u=x^2y^3+e^{x+y},
\]
find all first and second partial derivatives.

First:
\[
u_x=2xy^3+e^{x+y},
\qquad
u_y=3x^2y^2+e^{x+y}.
\]
Second:
\[
u_{xx}=2y^3+e^{x+y},
\]
\[
u_{xy}=6xy^2+e^{x+y},
\]
\[
u_{yx}=6xy^2+e^{x+y},
\]
\[
u_{yy}=6x^2y+e^{x+y}.
\]

\subsection*{Problem 2}
Find \(f_x(0,0)\) and \(f_y(0,0)\) directly for
\[
f(x,y)=
\begin{cases}
\dfrac{x^2y}{x^4+y^2},&(x,y)\ne(0,0),\\
0,&(0,0).
\end{cases}
\]
Since \(f(h,0)=0\),
\[
f_x(0,0)
=
\lim_{h\to0}\frac{0-0}{h}=0.
\]
Similarly \(f(0,k)=0\), so
\[
f_y(0,0)=0.
\]

\section{Graded Practice Problems}
\begin{enumerate}
\item Find all first partial derivatives of
\[
x^3+3x^2y+y^3.
\]
\item Find all second partial derivatives of \(e^{x+y}\).
\item Find \(f_x(1,2)\) and \(f_y(1,2)\) for \(x^2y+xy^2\).
\item Compute \(f_x(0,0),f_y(0,0)\) from first principles for
\(x^2+y^2\).
\item Give an example where both first partial derivatives exist but
the function is not continuous.
\end{enumerate}

% ============================================================
\chapter{Differentiability, Young's Theorem and Schwarz's Theorem}
% ============================================================

\section*{Learning Objectives}
Students should be able to:
\begin{itemize}
\item state the definition of differentiability in two variables;
\item interpret differentiability as linear approximation;
\item prove that differentiability implies continuity;
\item use continuity of first partial derivatives as a sufficient
condition for differentiability;
\item state Young's and Schwarz's theorems;
\item apply equality of mixed partial derivatives correctly.
\end{itemize}

\section*{Why Partial Derivatives Are Not Enough}

Partial derivatives only test changes along the coordinate directions.
Differentiability requires one linear approximation that works for
\emph{all} sufficiently small changes \((h,k)\).

\section{Definition of Differentiability}

\begin{definition}
A function \(f(x,y)\) is differentiable at \((a,b)\) if
\[
f(a+h,b+k)-f(a,b)
=
f_x(a,b)h+f_y(a,b)k
+
o\!\left(\sqrt{h^2+k^2}\right).
\]
Equivalently,
\[
\lim_{(h,k)\to(0,0)}
\frac{
f(a+h,b+k)-f(a,b)
-f_x(a,b)h-f_y(a,b)k
}{
\sqrt{h^2+k^2}
}
=0.
\]
\end{definition}

\paragraph{Geometrical meaning.}
Near \((a,b)\),
\[
f(a+h,b+k)
\approx
f(a,b)+f_x(a,b)h+f_y(a,b)k.
\]
Thus the surface is locally approximated by a plane.

The tangent plane is
\[
z
=
f(a,b)
+
f_x(a,b)(x-a)
+
f_y(a,b)(y-b).
\]

\section{Differentiability Implies Partial Derivatives Exist}

\begin{theorem}
If \(f\) is differentiable at \((a,b)\), then both \(f_x(a,b)\) and
\(f_y(a,b)\) exist.
\end{theorem}

\begin{proof}
Suppose
\[
f(a+h,b+k)-f(a,b)
=
Ah+Bk+\eta(h,k)\sqrt{h^2+k^2},
\]
where
\[
\eta(h,k)\to0.
\]

Set \(k=0\):
\[
f(a+h,b)-f(a,b)
=
Ah+\eta(h,0)\abs{h}.
\]
For \(h\ne0\),
\[
\frac{f(a+h,b)-f(a,b)}{h}
=
A+\eta(h,0)\frac{\abs h}{h}.
\]
Since
\[
\left|
\eta(h,0)\frac{\abs h}{h}
\right|
=
\abs{\eta(h,0)}
\to0,
\]
we obtain
\[
f_x(a,b)=A.
\]

Setting \(h=0\) similarly gives
\[
f_y(a,b)=B.
\]
\end{proof}

\section{Every Differentiable Function Is Continuous}

\begin{theorem}
If \(f(x,y)\) is differentiable at \((a,b)\), then \(f\) is
continuous at \((a,b)\).
\end{theorem}

\begin{proof}
To prove continuity we must show
\[
\lim_{(x,y)\to(a,b)}f(x,y)=f(a,b).
\]

Put
\[
x=a+h,\qquad y=b+k.
\]
Then
\[
(x,y)\to(a,b)
\quad\Longleftrightarrow\quad
(h,k)\to(0,0).
\]

Since \(f\) is differentiable,
\[
f(a+h,b+k)-f(a,b)
=
f_x(a,b)h
+
f_y(a,b)k
+
\eta(h,k)\sqrt{h^2+k^2},
\]
where
\[
\eta(h,k)\to0.
\]

Taking absolute values,
\[
\begin{aligned}
&
\abs{f(a+h,b+k)-f(a,b)}
\\
&\le
\abs{f_x(a,b)}\abs h
+
\abs{f_y(a,b)}\abs k
+
\abs{\eta(h,k)}\sqrt{h^2+k^2}.
\end{aligned}
\]

Since
\[
\abs h,\abs k\le\sqrt{h^2+k^2},
\]
we get
\[
\begin{aligned}
&
\abs{f(a+h,b+k)-f(a,b)}
\\
&\le
\left(
\abs{f_x(a,b)}
+\abs{f_y(a,b)}
+\abs{\eta(h,k)}
\right)
\sqrt{h^2+k^2}.
\end{aligned}
\]

As \((h,k)\to(0,0)\), the right-hand side tends to \(0\).
Therefore
\[
f(a+h,b+k)-f(a,b)\to0.
\]
Hence
\[
\lim_{(h,k)\to(0,0)}f(a+h,b+k)=f(a,b).
\]
Equivalently,
\[
\lim_{(x,y)\to(a,b)}f(x,y)=f(a,b).
\]
Thus \(f\) is continuous at \((a,b)\).
\end{proof}

\section{A Standard Sufficient Condition}

\begin{theorem}
Suppose \(f_x\) and \(f_y\) exist in a neighbourhood of \((a,b)\)
and are continuous at \((a,b)\). Then \(f\) is differentiable at
\((a,b)\).
\end{theorem}

\begin{proof}
Write
\[
\Delta f=f(a+h,b+k)-f(a,b).
\]
Insert and subtract \(f(a,b+k)\):
\[
\begin{aligned}
\Delta f
={}&
[f(a+h,b+k)-f(a,b+k)]\\
&+[f(a,b+k)-f(a,b)].
\end{aligned}
\]

Apply the one-variable Mean Value Theorem to the first difference as a
function of \(x\). For some \(\theta\in(0,1)\),
\[
f(a+h,b+k)-f(a,b+k)
=
h f_x(a+\theta h,b+k).
\]

Apply the Mean Value Theorem to the second difference as a function of
\(y\). For some \(\lambda\in(0,1)\),
\[
f(a,b+k)-f(a,b)
=
k f_y(a,b+\lambda k).
\]

Therefore
\[
\Delta f
=
h f_x(a+\theta h,b+k)
+
k f_y(a,b+\lambda k).
\]

Subtract
\[
f_x(a,b)h+f_y(a,b)k.
\]
Then
\[
\begin{aligned}
&
\Delta f-f_x(a,b)h-f_y(a,b)k\\
={}&
h\bigl[
f_x(a+\theta h,b+k)-f_x(a,b)
\bigr]\\
&+
k\bigl[
f_y(a,b+\lambda k)-f_y(a,b)
\bigr].
\end{aligned}
\]

Let
\[
\rho=\sqrt{h^2+k^2}.
\]
Then
\[
\frac{\abs h}{\rho}\le1,
\qquad
\frac{\abs k}{\rho}\le1.
\]
Hence the absolute value of the differentiability quotient is bounded
by
\[
\abs{f_x(a+\theta h,b+k)-f_x(a,b)}
+
\abs{f_y(a,b+\lambda k)-f_y(a,b)}.
\]
By continuity of \(f_x,f_y\), this tends to \(0\). Therefore \(f\) is
differentiable at \((a,b)\).
\end{proof}

\section{Logical Relationships}

It is useful to remember:
\[
\text{continuous first partial derivatives near the point}
\Longrightarrow
\text{differentiable}
\]
and
\[
\text{differentiable}
\Longrightarrow
\text{continuous}.
\]

But
\[
\text{partial derivatives exist}
\not\Rightarrow
\text{continuous}
\]
and therefore do not imply differentiability.

\section{Young's Theorem}

\begin{theorem}[Young's theorem --- statement]
Suppose the mixed partial derivatives \(f_{xy}\) and \(f_{yx}\)
exist in a neighbourhood of \((a,b)\) and one of them is continuous
at \((a,b)\), under the usual hypotheses ensuring the first partial
derivatives are defined. Then
\[
f_{xy}(a,b)=f_{yx}(a,b).
\]
\end{theorem}

Only the statement and applications are required here.

\section{Schwarz's Theorem}

\begin{theorem}[Schwarz's theorem --- statement]
If the second-order partial derivatives are continuous in a
neighbourhood of \((a,b)\), then
\[
f_{xy}(a,b)=f_{yx}(a,b).
\]
\end{theorem}

Again, only the statement and applications are required.

\section{Application}

For
\[
f(x,y)=x^3y^2+e^{xy},
\]
we obtain
\[
f_{xy}
=
6x^2y+e^{xy}+xye^{xy},
\]
and
\[
f_{yx}
=
6x^2y+e^{xy}+xye^{xy}.
\]
Both are continuous, so Schwarz's theorem applies.

\section{Why the Hypotheses Matter}

Consider
\[
f(x,y)=
\begin{cases}
\dfrac{xy(x^2-y^2)}{x^2+y^2},&(x,y)\ne(0,0),\\[1ex]
0,&(x,y)=(0,0).
\end{cases}
\]

For \(y\ne0\),
\[
f_x(0,y)
=
\lim_{h\to0}
\frac{f(h,y)-f(0,y)}{h}
=-y.
\]
Therefore
\[
f_{xy}(0,0)
=
\lim_{y\to0}
\frac{-y-0}{y}
=-1.
\]

Similarly one finds
\[
f_y(x,0)=x,
\]
hence
\[
f_{yx}(0,0)=1.
\]

Thus
\[
f_{xy}(0,0)\ne f_{yx}(0,0),
\]
showing that mixed partials are not automatically equal.

\section{Detailed Solved Problems}

\subsection*{Problem 1. Prove differentiability}
Show that
\[
f(x,y)=x^2+xy+y^2
\]
is differentiable everywhere.

Its first partial derivatives are
\[
f_x=2x+y,
\qquad
f_y=x+2y.
\]
These are polynomials, hence continuous everywhere. Therefore by the
standard sufficient condition, \(f\) is differentiable at every
point of \(\R^2\).

\subsection*{Problem 2. A harmonic-function application}
Let
\[
u(x,y)=e^x\cos y.
\]
Then
\[
u_x=e^x\cos y,
\qquad
u_{xx}=e^x\cos y.
\]
Also
\[
u_y=-e^x\sin y,
\qquad
u_{yy}=-e^x\cos y.
\]
Therefore
\[
u_{xx}+u_{yy}=0.
\]

\section{Graded Practice Problems}
\begin{enumerate}
\item Prove directly that differentiability implies continuity.
\item Determine whether
\[
f(x,y)=
\begin{cases}
\dfrac{x^2y}{x^2+y^2},&(x,y)\ne(0,0),\\
0,&(0,0)
\end{cases}
\]
is differentiable at the origin.
\item Verify Schwarz's theorem for
\[
f(x,y)=x^2y^3+\sin(xy).
\]
\item Give an example showing that mixed partial derivatives can
differ when the continuity hypotheses fail.
\end{enumerate}

% ============================================================
\chapter{Derivatives of Composite Functions}
% ============================================================

\section*{Learning Objectives}
Students should be able to:
\begin{itemize}
\item identify direct and indirect dependence among variables;
\item use the multivariable chain rule;
\item differentiate when all variables depend on one parameter;
\item use dependency diagrams to avoid missing terms.
\end{itemize}

\section*{Historical Motivation}

Many physical quantities depend on intermediate variables.
Temperature may depend on position, while position depends on time.
The chain rule organizes these layers of dependence and tells us how
small changes propagate through a system.

\section{Two Intermediate Variables}

Suppose
\[
z=f(u,v),
\]
where
\[
u=u(x,y),
\qquad
v=v(x,y).
\]

\begin{theorem}[Chain rule]
If the functions involved are differentiable, then
\[
\pd{z}{x}
=
\pd{z}{u}\pd{u}{x}
+
\pd{z}{v}\pd{v}{x},
\]
and
\[
\pd{z}{y}
=
\pd{z}{u}\pd{u}{y}
+
\pd{z}{v}\pd{v}{y}.
\]
\end{theorem}

\section{Derivation from Differentials}

For differentiable \(z=f(u,v)\),
\[
dz=f_u\,du+f_v\,dv.
\]
Now
\[
du=u_x\,dx+u_y\,dy,
\]
\[
dv=v_x\,dx+v_y\,dy.
\]
Substitution gives
\[
\begin{aligned}
dz
&=
f_u(u_x\,dx+u_y\,dy)
+
f_v(v_x\,dx+v_y\,dy)\\
&=
(f_u u_x+f_v v_x)\,dx
+
(f_u u_y+f_v v_y)\,dy.
\end{aligned}
\]
But
\[
dz=z_x\,dx+z_y\,dy.
\]
Comparing coefficients,
\[
z_x=f_u u_x+f_v v_x,
\]
\[
z_y=f_u u_y+f_v v_y.
\]

\section{Dependency Diagram}

\begin{center}
\begin{tikzpicture}[
node distance=1.5cm and 1.9cm,
every node/.style={draw,rounded corners,minimum width=1cm},
>=Stealth
]
\node (x) at (0,1) {$x$};
\node (y) at (0,-1) {$y$};
\node (u) at (3,1) {$u$};
\node (v) at (3,-1) {$v$};
\node (z) at (6,0) {$z$};

\draw[->] (x)--(u);
\draw[->] (x)--(v);
\draw[->] (y)--(u);
\draw[->] (y)--(v);
\draw[->] (u)--(z);
\draw[->] (v)--(z);
\end{tikzpicture}
\end{center}

For a derivative with respect to \(x\), add the products of
derivatives along every path from \(x\) to \(z\).

\section{Detailed Example}

Let
\[
z=u^2+v^2,
\qquad
u=x+y,
\qquad
v=x-y.
\]
First,
\[
z_u=2u,\qquad z_v=2v.
\]
Since
\[
u_x=1,\qquad v_x=1,
\]
we have
\[
z_x=2u+2v.
\]
Substituting \(u=x+y,v=x-y\),
\[
z_x=4x.
\]

Likewise
\[
u_y=1,\qquad v_y=-1,
\]
so
\[
z_y=2u-2v=4y.
\]

\section{One-Parameter Chain Rule}

If
\[
z=f(x,y),
\qquad
x=x(t),\quad y=y(t),
\]
then
\[
\frac{dz}{dt}
=
f_x\frac{dx}{dt}
+
f_y\frac{dy}{dt}.
\]

\begin{example}
Let
\[
z=x^2y,
\qquad
x=t^2,
\qquad
y=\sin t.
\]
Then
\[
z_x=2xy,\qquad z_y=x^2,
\]
and
\[
\frac{dx}{dt}=2t,
\qquad
\frac{dy}{dt}=\cos t.
\]
Therefore
\[
\begin{aligned}
\frac{dz}{dt}
&=
(2xy)(2t)+x^2\cos t\\
&=
4t^3\sin t+t^4\cos t.
\end{aligned}
\]

As a check, direct substitution gives
\[
z=t^4\sin t,
\]
whose derivative is the same expression.
\end{example}

\section{Detailed Solved Problems}

\subsection*{Problem 1}
Let
\[
z=e^{uv},
\qquad
u=x+y,
\qquad
v=x-y.
\]
Then
\[
z_u=ve^{uv},
\qquad
z_v=ue^{uv}.
\]
Thus
\[
z_x
=
ve^{uv}+ue^{uv}
=
(u+v)e^{uv}
=
2xe^{uv}.
\]
Similarly,
\[
z_y
=
ve^{uv}-ue^{uv}
=
(v-u)e^{uv}
=
-2ye^{uv}.
\]

\subsection*{Problem 2}
Let
\[
z=\log(u^2+v^2),
\quad
u=x^2-y^2,
\quad
v=2xy.
\]
Then
\[
z_u=\frac{2u}{u^2+v^2},
\qquad
z_v=\frac{2v}{u^2+v^2}.
\]
Also
\[
u_x=2x,\quad v_x=2y,
\qquad
u_y=-2y,\quad v_y=2x.
\]
Hence
\[
z_x
=
\frac{4xu+4yv}{u^2+v^2},
\]
and
\[
z_y
=
\frac{-4yu+4xv}{u^2+v^2}.
\]

\section{Graded Practice Problems}
\begin{enumerate}
\item If
\[
z=u^3+v^3,\quad u=x+y,\quad v=x-y,
\]
find \(z_x,z_y\).
\item If
\[
z=e^{xy},\quad x=t^2,\quad y=\cos t,
\]
find \(dz/dt\).
\item Draw a dependency diagram and derive the chain rule when
\(z=f(u,v,w)\) and each of \(u,v,w\) depends on \(x,y\).
\end{enumerate}

% ============================================================
\chapter{Change of Variables and Jacobians}
% ============================================================

\section*{Learning Objectives}
Students should be able to:
\begin{itemize}
\item understand why variables are changed;
\item calculate a Jacobian determinant;
\item derive transformation formulae for first derivatives;
\item handle a linear transformation;
\item transform derivatives between Cartesian and polar coordinates.
\end{itemize}

\section*{Historical Motivation}

Changing coordinates can reveal the geometry hidden in an algebraic
problem. Circular symmetry suggests polar coordinates, while other
geometries suggest other transformations. Jacobi's determinant gives a
systematic way to describe how variables change together.

\section{Change of Independent Variables}

Suppose
\[
u=u(x,y),
\qquad
v=v(x,y).
\]
If the transformation is invertible, we may also write
\[
x=x(u,v),
\qquad
y=y(u,v).
\]

A function
\[
z=f(x,y)
\]
may then be regarded as
\[
z=F(u,v).
\]

\section{The Jacobian}

\begin{definition}
The Jacobian of \(u,v\) with respect to \(x,y\) is
\[
\frac{\partial(u,v)}{\partial(x,y)}
=
\begin{vmatrix}
u_x&u_y\\
v_x&v_y
\end{vmatrix}
=
u_xv_y-u_yv_x.
\]
\end{definition}

A nonzero Jacobian is the elementary signal that the transformation is
locally invertible.

\section{A Linear Transformation}

Let
\[
u=x+y,
\qquad
v=x-y.
\]
Then
\[
u_x=1,\quad u_y=1,
\qquad
v_x=1,\quad v_y=-1.
\]
Therefore
\[
\frac{\partial(u,v)}{\partial(x,y)}
=
\begin{vmatrix}
1&1\\
1&-1
\end{vmatrix}
=-2.
\]

Solving the equations,
\[
x=\frac{u+v}{2},
\qquad
y=\frac{u-v}{2}.
\]

\section{Transformation of Differential Operators}

If \(z=F(u,v)\), then
\[
z_x=F_u u_x+F_v v_x.
\]
Symbolically,
\[
\frac{\partial}{\partial x}
=
u_x\frac{\partial}{\partial u}
+
v_x\frac{\partial}{\partial v}.
\]
Similarly,
\[
\frac{\partial}{\partial y}
=
u_y\frac{\partial}{\partial u}
+
v_y\frac{\partial}{\partial v}.
\]

For
\[
u=x+y,\qquad v=x-y,
\]
we obtain
\[
\frac{\partial}{\partial x}
=
\frac{\partial}{\partial u}
+
\frac{\partial}{\partial v},
\]
and
\[
\frac{\partial}{\partial y}
=
\frac{\partial}{\partial u}
-
\frac{\partial}{\partial v}.
\]

\section{A Second-Order Application}

Using the preceding formulas,
\[
z_{xx}
=
z_{uu}+2z_{uv}+z_{vv},
\]
and
\[
z_{yy}
=
z_{uu}-2z_{uv}+z_{vv}.
\]
Subtracting,
\[
z_{xx}-z_{yy}=4z_{uv}.
\]

\section{Reciprocal Jacobian}

If the transformation is invertible and
\[
J=
\frac{\partial(u,v)}{\partial(x,y)}
\ne0,
\]
then
\[
\frac{\partial(x,y)}{\partial(u,v)}
=
\frac1J.
\]

For the linear transformation above,
\[
\frac{\partial(x,y)}{\partial(u,v)}
=
-\frac12,
\]
which is indeed the reciprocal of \(-2\).

\section{Polar Coordinates}

Let
\[
x=r\cos\theta,
\qquad
y=r\sin\theta.
\]
If
\[
F(r,\theta)=f(x,y),
\]
then by the chain rule
\[
F_r
=
f_x\cos\theta+f_y\sin\theta.
\]

Also,
\[
x_\theta=-r\sin\theta,
\qquad
y_\theta=r\cos\theta,
\]
so
\[
F_\theta
=
-rf_x\sin\theta
+
rf_y\cos\theta.
\]

Solving these two equations for \(f_x,f_y\),
\[
f_x
=
F_r\cos\theta
-
\frac{\sin\theta}{r}F_\theta,
\]
\[
f_y
=
F_r\sin\theta
+
\frac{\cos\theta}{r}F_\theta.
\]

\section{Jacobian of Polar Coordinates}

For
\[
x=r\cos\theta,\qquad y=r\sin\theta,
\]
\[
x_r=\cos\theta,
\quad
x_\theta=-r\sin\theta,
\]
\[
y_r=\sin\theta,
\quad
y_\theta=r\cos\theta.
\]
Thus
\[
\begin{aligned}
\frac{\partial(x,y)}{\partial(r,\theta)}
&=
\begin{vmatrix}
\cos\theta&-r\sin\theta\\
\sin\theta&r\cos\theta
\end{vmatrix}\\
&=
r\cos^2\theta+r\sin^2\theta\\
&=r.
\end{aligned}
\]

\section{Detailed Solved Problems}

\subsection*{Problem 1. A quadratic transformation}
For
\[
u=x^2-y^2,
\qquad
v=2xy,
\]
find the Jacobian.

We have
\[
u_x=2x,\quad u_y=-2y,
\]
\[
v_x=2y,\quad v_y=2x.
\]
Therefore
\[
\begin{aligned}
\frac{\partial(u,v)}{\partial(x,y)}
&=
\begin{vmatrix}
2x&-2y\\
2y&2x
\end{vmatrix}\\
&=
4x^2+4y^2\\
&=
4(x^2+y^2).
\end{aligned}
\]

\subsection*{Problem 2. Transform a second-order expression}
Let
\[
u=x+y,\qquad v=x-y.
\]
Then
\[
\partial_x=\partial_u+\partial_v,
\qquad
\partial_y=\partial_u-\partial_v.
\]
Hence
\[
z_{xx}
=
z_{uu}+2z_{uv}+z_{vv},
\]
\[
z_{yy}
=
z_{uu}-2z_{uv}+z_{vv}.
\]
Therefore
\[
z_{xx}-z_{yy}=4z_{uv}.
\]

\section{Graded Practice Problems}
\begin{enumerate}
\item Find
\[
\frac{\partial(u,v)}{\partial(x,y)}
\]
for
\[
u=3x+2y,\qquad v=x-y.
\]
\item If
\[
u=xy,\qquad v=x/y,
\]
find the Jacobian.
\item Verify the reciprocal Jacobian relation for
\[
u=x+y,\qquad v=x-y.
\]
\item Derive the polar formulas for \(f_x,f_y\).
\item Find the Jacobian of
\[
x=r\cos\theta,\qquad y=r\sin\theta.
\]
\end{enumerate}

% ============================================================
\chapter{Homogeneous Functions and Euler's Theorem}
% ============================================================

\section*{Learning Objectives}
Students should be able to:
\begin{itemize}
\item identify homogeneous functions and their degree;
\item state and prove Euler's theorem for two variables;
\item derive first- and second-order consequences;
\item apply Euler's theorem to composite functions efficiently.
\end{itemize}

\section*{Historical Motivation}

Homogeneous functions encode scaling. If all independent variables are
multiplied by the same factor, the output changes by a predictable
power of that factor. Such functions appear throughout geometry,
mechanics, thermodynamics and mathematical economics.

\section{Definition}

\begin{definition}
A function \(u=f(x,y)\) is homogeneous of degree \(n\) if
\[
f(tx,ty)=t^n f(x,y)
\]
for all suitable \(t>0\).
\end{definition}

\begin{example}
For
\[
f(x,y)=x^2+xy+y^2,
\]
\[
f(tx,ty)
=
t^2(x^2+xy+y^2).
\]
Hence \(f\) is homogeneous of degree \(2\).
\end{example}

\begin{example}
For
\[
f(x,y)=\frac{x+y}{x-y},
\]
\[
f(tx,ty)=f(x,y).
\]
Hence the degree is \(0\).
\end{example}

\section{Euler's Theorem}

\begin{theorem}[Euler's theorem for two variables]
Let
\[
u=f(x,y)
\]
be differentiable and homogeneous of degree \(n\). Then
\[
x\pd{u}{x}
+
y\pd{u}{y}
=
nu.
\]
\end{theorem}

\begin{proof}
Since \(u=f(x,y)\) is homogeneous of degree \(n\),
\[
f(tx,ty)=t^n f(x,y).
\]

Treat \(x,y\) as constants and differentiate both sides with respect
to \(t\).

By the chain rule,
\[
\begin{aligned}
\frac{d}{dt}f(tx,ty)
&=
f_x(tx,ty)\frac{d(tx)}{dt}
+
f_y(tx,ty)\frac{d(ty)}{dt}\\
&=
x f_x(tx,ty)
+
y f_y(tx,ty).
\end{aligned}
\]

On the right,
\[
\frac{d}{dt}
[t^n f(x,y)]
=
nt^{n-1}f(x,y).
\]

Therefore
\[
x f_x(tx,ty)
+
y f_y(tx,ty)
=
nt^{n-1}f(x,y).
\]

Put \(t=1\):
\[
x f_x(x,y)+y f_y(x,y)=nf(x,y).
\]
Since \(u=f(x,y)\),
\[
xu_x+yu_y=nu.
\]
\end{proof}

\section{First Differential Consequences}

Starting with
\[
xu_x+yu_y=nu,
\]
differentiate with respect to \(x\):
\[
u_x+xu_{xx}+yu_{yx}=nu_x.
\]
Assuming equality of the mixed derivatives under the appropriate
hypotheses,
\[
xu_{xx}+yu_{xy}=(n-1)u_x.
\]

Similarly, differentiation with respect to \(y\) gives
\[
xu_{xy}+yu_{yy}=(n-1)u_y.
\]

\section{Second-Order Euler Identity}

Multiply
\[
xu_{xx}+yu_{xy}=(n-1)u_x
\]
by \(x\):
\[
x^2u_{xx}+xyu_{xy}=(n-1)xu_x.
\]

Multiply
\[
xu_{xy}+yu_{yy}=(n-1)u_y
\]
by \(y\):
\[
xyu_{xy}+y^2u_{yy}=(n-1)yu_y.
\]

Add:
\[
x^2u_{xx}
+
2xyu_{xy}
+
y^2u_{yy}
=
(n-1)(xu_x+yu_y).
\]

By Euler's theorem,
\[
xu_x+yu_y=nu.
\]
Hence
\[
x^2u_{xx}
+
2xyu_{xy}
+
y^2u_{yy}
=
n(n-1)u.
\]

\section{Composite Form of Euler's Theorem}

Suppose \(v(x,y)\) is homogeneous of degree \(n\) and
\[
u=\phi(v).
\]
Then
\[
u_x=\phi'(v)v_x,
\qquad
u_y=\phi'(v)v_y.
\]
Therefore
\[
\begin{aligned}
xu_x+yu_y
&=
\phi'(v)(xv_x+yv_y)\\
&=
nv\phi'(v).
\end{aligned}
\]

This simple observation often saves a large amount of direct
differentiation.

\section{Detailed Solved Problems}

\subsection*{Problem 1 (2025). Euler's theorem}
Let \(u=f(x,y)\) be homogeneous of degree \(n\). Show that
\[
xu_x+yu_y=nu.
\]

Since
\[
f(tx,ty)=t^n f(x,y),
\]
differentiate with respect to \(t\):
\[
x f_x(tx,ty)+y f_y(tx,ty)
=
nt^{n-1}f(x,y).
\]
Putting \(t=1\) gives
\[
xu_x+yu_y=nu.
\]

\subsection*{Problem 2 (2024). A trigonometric application}
Let
\[
u
=
\tan^{-1}
\left(
\frac{x^3+y^3}{x+y}
\right).
\]
Show that
\[
xu_x+yu_y=\sin2u.
\]

Set
\[
v=\frac{x^3+y^3}{x+y}.
\]
Since
\[
x^3+y^3=(x+y)(x^2-xy+y^2),
\]
we have
\[
v=x^2-xy+y^2,
\]
where the original expression is defined. Thus \(v\) is homogeneous
of degree \(2\).

Euler's theorem gives
\[
xv_x+yv_y=2v.
\]

Now
\[
u=\tan^{-1}v,
\]
so
\[
u_x=\frac{v_x}{1+v^2},
\qquad
u_y=\frac{v_y}{1+v^2}.
\]
Therefore
\[
\begin{aligned}
xu_x+yu_y
&=
\frac{xv_x+yv_y}{1+v^2}\\
&=
\frac{2v}{1+v^2}.
\end{aligned}
\]

Since
\[
v=\tan u,
\]
\[
\begin{aligned}
xu_x+yu_y
&=
\frac{2\tan u}{1+\tan^2u}\\
&=
2\tan u\cos^2u\\
&=
2\sin u\cos u\\
&=
\sin2u.
\end{aligned}
\]

\subsection*{Problem 3 (2023). A differentiated Euler identity}
Suppose \(u(x,y)\) is homogeneous of degree \(n\). Simplify
\[
xu_{xy}+yu_{yy}.
\]

Euler's theorem gives
\[
xu_x+yu_y=nu.
\]
Differentiate with respect to \(y\):
\[
xu_{xy}+u_y+yu_{yy}=nu_y.
\]
Hence
\[
xu_{xy}+yu_{yy}=(n-1)u_y.
\]

\subsection*{Problem 4. Logarithm of a homogeneous function}
Suppose \(v(x,y)\) is homogeneous of degree \(n\), and
\[
u=\log v.
\]
Then
\[
u_x=\frac{v_x}{v},
\qquad
u_y=\frac{v_y}{v}.
\]
Thus
\[
xu_x+yu_y
=
\frac{xv_x+yv_y}{v}
=
\frac{nv}{v}
=
n.
\]

\section{Graded Practice Problems}

\subsection*{Level I}
\begin{enumerate}
\item Determine the degree of
\[
x^3+x^2y+xy^2+y^3.
\]
\item Determine the degree of
\[
\frac{x^2+y^2}{x+y}.
\]
\item Verify Euler's theorem for
\[
u=x^2+xy+y^2.
\]
\end{enumerate}

\subsection*{Level II}
\begin{enumerate}
\item Verify Euler's theorem for
\[
u=(x^2+y^2)^{3/2}.
\]
\item If \(u\) is homogeneous of degree \(n\), prove
\[
xu_{xx}+yu_{xy}=(n-1)u_x.
\]
\item Derive
\[
x^2u_{xx}+2xyu_{xy}+y^2u_{yy}=n(n-1)u.
\]
\end{enumerate}

\subsection*{Level III}
\begin{enumerate}
\item If \(v\) is homogeneous of degree \(n\) and \(u=\log v\),
show that \(xu_x+yu_y=n\).
\item If \(v\) is homogeneous of degree \(n\) and \(u=e^v\), show
that
\[
xu_x+yu_y=nve^v.
\]
\item If \(v\) is homogeneous of degree \(n\) and
\(u=\sin^{-1}v\), express \(xu_x+yu_y\) in terms of \(u\).
\end{enumerate}

% ============================================================
\chapter{Integrated Review and Problem-Solving Guide}
% ============================================================

\section{How to Approach a Two-Variable Limit}

When asked to examine
\[
\lim_{(x,y)\to(a,b)}f(x,y),
\]
use the following sequence.

\begin{enumerate}
\item Try direct substitution.
\item If an indeterminate form occurs, simplify algebraically.
\item If the point is the origin and \(x^2+y^2\) is present, consider
polar coordinates.
\item To prove nonexistence, test simple paths:
\[
x=0,\qquad y=0,\qquad y=mx.
\]
\item If these paths agree, do not conclude existence. Consider curved
paths such as \(y=x^2\) when suggested by the algebra.
\item To prove existence, use an estimate, the Squeeze Theorem, polar
coordinates, or the \(\eps\)-\(\delta\) definition.
\end{enumerate}

\section{How to Test Differentiability}

At \((a,b)\):

\begin{enumerate}
\item Check continuity first. A discontinuous function cannot be
differentiable.
\item Find \(f_x(a,b)\) and \(f_y(a,b)\).
\item If \(f_x,f_y\) exist in a neighbourhood and are continuous at
the point, differentiability follows.
\item Otherwise test
\[
\frac{
f(a+h,b+k)-f(a,b)-f_x(a,b)h-f_y(a,b)k
}{
\sqrt{h^2+k^2}
}
\to0.
\]
\end{enumerate}

\section{Common Misconceptions}

\begin{enumerate}
\item Equal values along the coordinate axes do not prove a joint
limit exists.
\item Agreement along all straight lines still need not prove a limit
exists.
\item Equal repeated limits do not guarantee the joint limit.
\item Existence of \(f_x,f_y\) does not imply continuity.
\item Existence of \(f_x,f_y\) does not imply differentiability.
\item Mixed partial derivatives are not automatically equal without
suitable hypotheses.
\item Euler's theorem applies only after homogeneity and its degree
have been identified correctly.
\item In a chain-rule problem, every path of dependence contributes a
term.
\end{enumerate}

\section{Quick Reference}

\subsection*{Joint limit}
\[
\lim_{(x,y)\to(a,b)}f(x,y)=L
\]
means: for every \(\eps>0\), there exists \(\delta>0\) such that
\[
0<\sqrt{(x-a)^2+(y-b)^2}<\delta
\]
implies
\[
\abs{f(x,y)-L}<\eps.
\]

\subsection*{Continuity}
\[
f\text{ continuous at }(a,b)
\iff
\lim_{(x,y)\to(a,b)}f(x,y)=f(a,b).
\]

\subsection*{Partial derivatives}
\[
f_x(a,b)
=
\lim_{h\to0}
\frac{f(a+h,b)-f(a,b)}{h},
\]
\[
f_y(a,b)
=
\lim_{k\to0}
\frac{f(a,b+k)-f(a,b)}{k}.
\]

\subsection*{Differentiability}
\[
f(a+h,b+k)-f(a,b)
=
f_x(a,b)h+f_y(a,b)k
+
o(\sqrt{h^2+k^2}).
\]

\subsection*{Tangent plane}
\[
z=f(a,b)+f_x(a,b)(x-a)+f_y(a,b)(y-b).
\]

\subsection*{Chain rule}
If
\[
z=f(u,v),\quad u=u(x,y),\quad v=v(x,y),
\]
then
\[
z_x=f_u u_x+f_v v_x,
\qquad
z_y=f_u u_y+f_v v_y.
\]

\subsection*{Jacobian}
\[
\frac{\partial(u,v)}{\partial(x,y)}
=
\begin{vmatrix}
u_x&u_y\\
v_x&v_y
\end{vmatrix}.
\]

\subsection*{Euler's theorem}
If \(u\) is homogeneous of degree \(n\),
\[
xu_x+yu_y=nu.
\]
Also,
\[
xu_{xx}+yu_{xy}=(n-1)u_x,
\]
\[
xu_{xy}+yu_{yy}=(n-1)u_y,
\]
and
\[
x^2u_{xx}+2xyu_{xy}+y^2u_{yy}=n(n-1)u.
\]

\section{Comprehensive Review Problems}

\begin{enumerate}
\item Examine
\[
\lim_{(x,y)\to(0,0)}
\frac{x^2y}{x^4+y^2}.
\]

\item Determine both repeated limits of
\[
\frac{x-y}{x+y}.
\]

\item Find all first- and second-order partial derivatives of
\[
u=e^{xy}\sin(x+y).
\]

\item Determine whether
\[
f(x,y)=
\begin{cases}
\dfrac{x^2y}{x^2+y^2},&(x,y)\ne(0,0),\\
0,&(0,0)
\end{cases}
\]
is continuous at the origin.

\item Explain why existence of both first partial derivatives at a
point does not guarantee differentiability.

\item Let
\[
z=u^2v,\quad u=x+y,\quad v=x-y.
\]
Find \(z_x,z_y\).

\item For
\[
u=x^2-y^2,\qquad v=2xy,
\]
calculate
\[
\frac{\partial(u,v)}{\partial(x,y)}.
\]

\item If \(u\) is homogeneous of degree \(5\), simplify
\[
xu_x+yu_y.
\]

\item If \(u\) is homogeneous of degree \(n\), simplify
\[
xu_{xy}+yu_{yy}.
\]

\item If
\[
v=x^2+xy+y^2,
\qquad
u=\log v,
\]
find \(xu_x+yu_y\) without separately expanding both partial
derivatives.
\end{enumerate}

\backmatter

\chapter*{Final Perspective}
\addcontentsline{toc}{chapter}{Final Perspective}

The essential change from one-variable to multivariable calculus is
geometrical. A point on a line has only two basic directions of
approach, while a point in the plane can be approached along infinitely
many curves. A one-variable derivative gives one local slope; a
two-variable function has several directional rates of change.

Partial derivatives isolate coordinate directions, while
differentiability brings all small changes together through a single
linear approximation. The chain rule explains how changes propagate
through intermediate variables. Jacobians organize changes of
coordinates, and Euler's theorem converts a scaling law into a
differential identity.

These ideas are foundational for vector calculus, partial differential
equations, differential geometry, mathematical physics and applied
mathematics.

\end{document}
